\subsection{Prompt Templates}\label{sec:app-prompts}
\textbf{Prompt Configuration for Question Answering and Reasoning.}
We would like to highlight that our evaluation protocol encourages the diffusion model to generate the output tokens with the sense of reasoning for \textbf{both} math solving and basic common QA tasks~(e.g., PiQA). As a result, the model achieves higher accuracy compared to the reported performance in \cite{dream2025}. To solidify the fact that we are not heavily relying on prompt engineering, we release the prompt templates as follows:



% Each prompt in the following is used verbatim. Replace ``$<$sentence$>$'', ``$<$question$>$'', ``$<$problem$>$'', or ``$<$instruction$>$'' with the actual input. All ``**'' markers are preserved exactly as in the original implementation.

\medskip
\begin{tcolorbox}[colback=gray!5, colframe=gray!50, boxrule=0.5pt, arc=2pt]
\textbf{ARC (-C/E):}\\

You are a careful reasoning assistant.  
A commonsense sentence is shown with a blank and **four** Answers that could fill the blank.  
Think step by step about which option makes the sentence most sensible, then decide.\\

\#\#\# Sentence:  
$<$sentence$>$ \\

\#\#\# What you should do:  \\
1. Briefly explain your reasoning.  \\
2. On a **new line**, write exactly:  
   "The final answer is [answer]"  
   where **[answer]** is either **Answer1**, **Answer2**, **Answer3**, or **Answer4**.  
Do **not** output anything after that line.
\end{tcolorbox}

\medskip
\begin{tcolorbox}[colback=gray!5, colframe=gray!50, boxrule=0.5pt, arc=2pt]
\textbf{GPQA:}\\

You are a careful reasoning assistant.  
A question is shown with **four** possible answers.  
Think step by step about which option is correct, then decide. \\

\#\#\# Question: $<$question$>$ \\

\#\#\# In your response:  \\
1. First reason and explain your reasoning briefly before providing the final answer to the question.  \\
2. Then at the end of your explanation, on a **new line**, write exactly:  
   "The final answer is [answer]"  
   where **[answer]** is one of **Answer1**, **Answer2**, **Answer3**, or **Answer4**.  
Do **not** output anything after that line.
\end{tcolorbox}

\medskip
\begin{tcolorbox}[colback=gray!5, colframe=gray!50, boxrule=0.5pt, arc=2pt]
\textbf{GSM8K:}\\

Given the following problem, reason and give a final answer to the problem. \\

Problem:  
$<$problem$>$ \\

Your response should end with  
"The final answer is [answer]"  
where [answer] is just the final number to the problem.
\end{tcolorbox}

\medskip
\begin{tcolorbox}[colback=gray!5, colframe=gray!50, boxrule=0.5pt, arc=2pt]
\textbf{HellaSwag:} \\

You are a careful reasoning assistant.  
A commonsense sentence is shown with a blank and **four** Endings that could fill the blank.  
Think step by step about which option makes the sentence most sensible, then decide. \\

\#\#\# Sentence:  
$<$sentence$>$ \\

\#\#\# In your response:  \\
1. First reason and explain your reasoning briefly before answering this question.  \\
2. Then at the end of your explanation, on a **new line**, write exactly:  
   "The final answer is [answer]"  
   where **[answer]** is either **Ending1**, **Ending2**, **Ending3**, or **Ending4**.  
Do **not** output anything after that line.
\end{tcolorbox}

% \medskip
% \begin{tcolorbox}[colback=gray!5, colframe=gray!50, boxrule=0.5pt, arc=2pt]
% \textbf{OBQA:} \\

% You are a careful reasoning assistant.  
% A commonsense sentence is shown with a blank and **four** Answers that could fill the blank.  
% Think step by step about which option makes the sentence most sensible, then decide. \\


% \#\#\# Sentence:  
% $<$sentence$>$ \\

% \#\#\# What you should do:  


% 1. Briefly explain your reasoning.  \\
% 2. On a **new line**, write exactly:  
%    "The final answer is [answer]"  
%    where **[answer]** is either **Answer1**, **Answer2**, **Answer3**, or **Answer4**.  
% Do **not** output anything after that line.
% \end{tcolorbox}

\medskip
\begin{tcolorbox}[colback=gray!5, colframe=gray!50, boxrule=0.5pt, arc=2pt]
\textbf{PIQA:} \\

Below is an instruction that describes a task. Write a response that appropriately completes the request. \\

\#\#\# Instruction:  
$<$instruction$>$ \\

Your response should end with  
"The final answer is [answer]"  
where [answer] is the response to the problem.
\end{tcolorbox}

\medskip
\begin{tcolorbox}[colback=gray!5, colframe=gray!50, boxrule=0.5pt, arc=2pt]
\textbf{MMLU-PRO:} \\

The following are multiple choice questions (with answers) about a subject. Think step by step and then finish your answer with "the answer is (X)" where X is the correct letter choice.
\end{tcolorbox}
The prompt of MMLU-PRO is adopted from the official repository of MMLU-PRO benchmark. 

% \medskip
% \begin{tcolorbox}[colback=gray!5, colframe=gray!50, boxrule=0.5pt, arc=2pt]
% \textbf{Winogrande:} \\

% You are a careful reasoning assistant.  
% A commonsense sentence is shown with a blank and **two** options that could fill the blank.  
% Think step by step about which option makes the sentence most sensible, then decide. \\

% \#\#\# Sentence:  
% $<$sentence$>$ \\

% \#\#\# What you should do:  \\
% 1. Briefly explain your reasoning.  
% 2. On a **new line**, write exactly:  \\
%    "The final answer is [answer]"  
%    where **[answer]** must be either **option1** or **option2**.  

% Do **not** output anything after that line.
% \end{tcolorbox}


% \section{Prompt Templates}\label{sec:app-prompts}

% Each prompt below was used verbatim. Replace `$<$sentence$>$`, `$<$question$>$`, `$<$problem$>$`, or `$<$instruction$>$` with the actual input. All `**` markers are preserved exactly as in the original implementation.

% \medskip
% \begin{tcolorbox}[colback=gray!5, colframe=gray!50, boxrule=0.5pt, arc=2pt]
% ARCC:

% You are a careful reasoning assistant.  
% A commonsense sentence is shown with a blank and **four** Answers that could fill the blank.  
% Think step by step about which option makes the sentence most sensible, then decide.

% \#\#\# Sentence:  
% $<$sentence$>$

% \#\#\# What you should do:  
% 1. Briefly explain your reasoning.  
% 2. On a **new line**, write exactly:  
%    "The final answer is [answer]"  
%    where **[answer]** is either **Answer1**, **Answer2**, **Answer3**, or **Answer4**.  
% Do **not** output anything after that line.
% \end{tcolorbox}

% \medskip
% \begin{tcolorbox}[colback=gray!5, colframe=gray!50, boxrule=0.5pt, arc=2pt]
% GPQA:

% You are a careful reasoning assistant.  
% A question is shown with **four** possible answers.  
% Think step by step about which option is correct, then decide.

% $<$question$>$

% \#\#\# In your response:  
% 1. First reason and explain your reasoning briefly before providing the final answer to the question.  
% 2. Then at the end of your explanation, on a **new line**, write exactly:  
%    "The final answer is [answer]"  
%    where **[answer]** is one of **Answer1**, **Answer2**, **Answer3**, or **Answer4**.  
% Do **not** output anything after that line.
% \end{tcolorbox}

% \medskip
% \begin{tcolorbox}[colback=gray!5, colframe=gray!50, boxrule=0.5pt, arc=2pt]
% GSM8K:

% Given the following problem, reason and give a final answer to the problem.

% Problem:  
% $<$problem$>$

% Your response should end with  
% "The final answer is [answer]"  
% where [answer] is just the final number to the problem.
% \end{tcolorbox}

% \medskip
% \begin{tcolorbox}[colback=gray!5, colframe=gray!50, boxrule=0.5pt, arc=2pt]
% HellaSwag:

% You are a careful reasoning assistant.  
% A commonsense sentence is shown with a blank and **four** Endings that could fill the blank.  
% Think step by step about which option makes the sentence most sensible, then decide.

% \#\#\# Sentence:  
% $<$sentence$>$

% \#\#\# In your response:  
% 1. First reason and explain your reasoning briefly before answering this question.  
% 2. Then at the end of your explanation, on a **new line**, write exactly:  
%    "The final answer is [answer]"  
%    where **[answer]** is either **Ending1**, **Ending2**, **Ending3**, or **Ending4**.  
% Do **not** output anything after that line.
% \end{tcolorbox}

% \medskip
% \begin{tcolorbox}[colback=gray!5, colframe=gray!50, boxrule=0.5pt, arc=2pt]
% OBQA:

% You are a careful reasoning assistant.  
% A commonsense sentence is shown with a blank and **four** Answers that could fill the blank.  
% Think step by step about which option makes the sentence most sensible, then decide.

% \#\#\# Sentence:  
% $<$sentence$>$

% \#\#\# What you should do:  
% 1. Briefly explain your reasoning.  
% 2. On a **new line**, write exactly:  
%    "The final answer is [answer]"  
%    where **[answer]** is either **Answer1**, **Answer2**, **Answer3**, or **Answer4**.  
% Do **not** output anything after that line.
% \end{tcolorbox}

% \medskip
% \begin{tcolorbox}[colback=gray!5, colframe=gray!50, boxrule=0.5pt, arc=2pt]
% PIQA:

% Below is an instruction that describes a task. Write a response that appropriately completes the request.

% \#\#\# Instruction:  
% $<$instruction$>$

% Your response should end with  
% "The final answer is [answer]"  
% where [answer] is the response to the problem.
% \end{tcolorbox}

% \medskip
% \begin{tcolorbox}[colback=gray!5, colframe=gray!50, boxrule=0.5pt, arc=2pt]
% Winogrande:

% You are a careful reasoning assistant.  
% A commonsense sentence is shown with a blank and **two** options that could fill the blank.  
% Think step by step about which option makes the sentence most sensible, then decide.

% \#\#\# Sentence:  
% $<$sentence$>$

% \#\#\# What you should do:  
% 1. Briefly explain your reasoning.  
% 2. On a **new line**, write exactly:  
%    "The final answer is [answer]"  
%    where **[answer]** must be either **option1** or **option2**.  
% Do **not** output anything after that line.
% \end{tcolorbox}
%  Sentence:  
% $<$sentence$>$

% \#\#\# What you should do:  
% 1. Briefly explain your reasoning.  
% 2. On a **new line**, write exactly:  
%    "The final answer is [answer]"  
%    where **[answer]** must be either **option1** or **option2**.  
% Do **not** output anything after that line.
% \end{tcolorbox}
